\documentclass[11pt,a4paper]{article}

\usepackage[utf8]{inputenc}
\usepackage[T1]{fontenc}
\usepackage{amsmath,amssymb,amsfonts}
\usepackage{booktabs}
\usepackage{array}
\usepackage{tabularx}
\usepackage{geometry}
\usepackage{hyperref}
\usepackage{xcolor}
\usepackage[protrusion=true,expansion=false]{microtype}

\geometry{margin=1in}

\hypersetup{
    colorlinks=true,
    linkcolor=blue!70!black,
    citecolor=green!50!black,
    urlcolor=blue!80!black
}

\title{\textbf{Simscape Multibody Golf Swing Modeling: Mechanics, Identifiability, and Diagnostics}}
\author{UpstreamDrift Modeling Reference}
\date{October 1, 2026 --- Source Revision \texttt{144e81188}}

\begin{document}

\maketitle

\begin{abstract}
This reference documents the Simscape Multibody golf swing model lineage from the legacy kinetic model to the contact-supported \texttt{GS3DX\_Human} variant under MATLAB R2025b (Home license: 1,000-block limit, 25-block diagnostic reserve). It details coordinate transformations, address waist origins, intrinsic Euler mappings, and power-preserving torque duals. We record parameter identifiability limits, iterative learning control with inertia-scaled gains, foot-pose Jacobian balance feedback, and closed-loop internal force spaces. Point-force contact diagnostics (21.863~N$\cdot$m RMS residual) and transient load spikes (365~BW and 22.6~BW) are evaluated, rejecting unqualified contact loads. Active red acceptance gates (\#11156, \#11160, \#11173) and operational qualification protocols are established.
\end{abstract}

\tableofcontents
\newpage

\section{Document Governance and Authority Notice}

\subsection{Canonical Manual Governance Relationship}
This document serves strictly as a dedicated research, experiment, and biomechanical modeling reference record for the exploratory Simscape Multibody matching continuation. In compliance with repository policy specified in \texttt{CLAUDE.md}, \texttt{AGENTS.md}, and \nolinkurl{scripts/config/design_manual_governance.json}:
\begin{enumerate}
    \item The sole editable scientific authority for the calculation-level engineering design manual is located at \texttt{manuals/upstreamdrift} (QMD source).
    \item This reference document is \textbf{a separate research and experiment record, not the engineering design manual}, and does not alter the canonical calculation registry.
    \item Generated outputs of the canonical manual remain non-editable artifacts governed by automated cryptographic provenance checks.
    \item This document is maintained independently within \nolinkurl{docs/research/simscape_matching_reference/} to record experimental findings, architectural derivations, diagnostic falsifications, and handoff criteria.
\end{enumerate}

\subsection{Branch Provenance and Baseline Commit Ledger}
Historical metrics are attributed to the October 1 turnover.\par
Turnover branch: \nolinkurl{feat/simscape-gs3dx-exploratory}.\par
Reviewed branch: \nolinkurl{feat/simscape-reference-20261001}.
\begin{itemize}
    \item \textbf{Branch Head Source Commit:} \texttt{144e81188dd7bb106f81d89b7a7330dd20cce511} (October 1, 2026).
    \item \textbf{Historical Handoff Baseline:} \texttt{2d5830d18}. This historical attribution is not a verified current remote baseline.
    \item \textbf{Active Draft PR:} \#11179 (exploratory GS3DX continuation).
    \item \textbf{Primary Governing Issues:} Epic \#10950 (children \#10951--\#10959, \#10985, \#10986, \#11011; active \#10979); Quality gates \#11156, \#11157, \#11158, \#11159, \#11160; Forward dynamics objective \#11173; Subject matching epic \#11161.
\end{itemize}

\subsection{Data Anonymization and Privacy Policy}
In strict adherence to fleet-wide data security standards, no proprietary athlete identities, personal health identifiers, commercial measurement hardware vendor names, or private absolute local filesystem paths are committed to version control:
\begin{itemize}
    \item All private motion capture directories are resolved at runtime via the environment variable \texttt{\$env:CAPTURE\_DATA\_DIR}.
    \item Motion capture datasets are referenced via standardized anonymized aliases:
    \begin{itemize}
        \item \texttt{capture-A}: Tour driver reference dataset; no population-average qualification is implied.
        \item \texttt{capture-B}: Secondary professional trial datasets.
        \item \texttt{capture-O}: Custom subject evaluation dataset (1.956~m height, 104.3~kg mass).
    \end{itemize}
\end{itemize}

\section{Coordinate Frames, Units, and Sign Conventions}

\subsection{Global World, Capture, and Address Coordinate Systems}
The simulation environment enforces standard SI metric units: lengths in meters (m), masses in kilograms (kg), forces in Newtons (N), torques in Newton-meters (N$\cdot$m), angles in radians (rad) for mathematical formulations, and degrees ($^\circ$) for tabular reporting.

\begin{itemize}
    \item \textbf{Simscape World Frame ($\mathcal{F}_W$):} Defined with $+Z_W$ vertical against gravity ($\mathbf{g} = [0, 0, -9.80665]^T$~m/s$^2$), $+X_W$ pointing forwards across the golfer's frontal plane at address (facing direction), and $+Y_W$ pointing laterally down the target line to the golfer's left.
    \item \textbf{Native Capture Coordinate Frame ($\mathcal{F}_C$):} Optical motion capture systems natively export trajectories with $+Y_C$ pointing upwards. The loader \texttt{gs3dx\_capture\_markers.m} performs an orthogonal mapping into the Simscape world frame:
    \begin{equation}
        \mathbf{p}_W = \begin{bmatrix} 1 & 0 & 0 \\ 0 & 0 & -1 \\ 0 & 1 & 0 \end{bmatrix} \mathbf{p}_C = \begin{bmatrix} x_C \\ -z_C \\ y_C \end{bmatrix}.
    \end{equation}
    \item \textbf{Address Frame ($\mathcal{F}_{\text{addr}}$) and Waist Origin:} In \texttt{gs3dx\_capture\_address\_transform.m}, the address origin $\mathbf{p}_{\text{origin}}$ is defined as the mean position of the \textbf{four waist markers} at frame 1 ($t = 0$):
    \begin{equation}
        \mathbf{p}_{\text{origin}} = \frac{1}{4} \left(\mathbf{p}_{\text{WaistLeft}} + \mathbf{p}_{\text{WaistRight}} + \mathbf{p}_{\text{WaistLBack}} + \mathbf{p}_{\text{WaistRBack}}\right)_{t=0}.
    \end{equation}
    The address origin is \textbf{not} the ankle midpoint or ground datum. Point positions transform to local address coordinates via the basis matrix $\mathbf{S} = [\mathbf{s}_{\text{facing}}, \mathbf{s}_{\text{target}}, \mathbf{s}_{\text{up}}]$:
    \begin{equation}
        \mathbf{p}_{\text{addr}} = \mathbf{S}^T (\mathbf{p}_W - \mathbf{p}_{\text{origin}}).
    \end{equation}
    Under this change of frame, forces transform purely by rotation:
    \begin{equation}
        \mathbf{f}_{\text{addr}} = \mathbf{S}^T \mathbf{f}_W,
    \end{equation}
    For a resultant moment $\mathbf{m}_W$ taken about the World origin,
    moments transform by incorporating the lever arm of the origin shift:
    \begin{equation}
        \mathbf{m}_{\text{addr}} = \mathbf{S}^T \mathbf{m}_W - (\mathbf{S}^T \mathbf{p}_{\text{origin}}) \times \mathbf{f}_{\text{addr}}.
    \end{equation}
    \item \textbf{Ground Datum vs. Address Origin:} The waist origin sits at elevated vertical position. The ground height belongs to the selected model workspace and experiment. A recorded run used $Z_W \approx -1.00027$~m and ankle height $0.0702$~m; these are not universal constants or the capture address origin. Record the actual floor transform for each replay.
\end{itemize}

\subsection{Joint Sign Conventions and Anatomical Definitions}
Joint rotations are parameterized as intrinsic transformations. For revolute and universal primitives, the anatomical angle $A_j$ is related to the raw Simscape generalized coordinate $q_j$ via:
\begin{equation}
    A_j = s_j \cdot \text{wrap}_{180}(q_j - q_{j,0}),
\end{equation}
where $s_j \in \{+1, -1\}$ is the structural joint direction sign, $q_{j,0}$ is the anatomical neutral angle, and $\text{wrap}_{180}(\theta) = \theta - 360^\circ \cdot \lfloor(\theta + 180^\circ)/360^\circ\rfloor$.

\begin{table}[!htb]
\centering
\footnotesize
\caption{Anatomical Joint Kinematic Signs, Neutrals, and Primary Physiological Boundaries}
\label{tab:joint_signs}
\begin{tabularx}{\textwidth}{p{2.6cm} >{\raggedright\arraybackslash}X c p{3.4cm} >{\raggedright\arraybackslash}X}
\toprule
\textbf{Joint Axis / Type} & \textbf{Positive Motion ($s_j$)} & \textbf{Neutral ($q_0$)} & \textbf{ROM Range} & \textbf{Primary Source} \\
\midrule
Neck Rx \newline \textit{Universal} & Trunk lateral bend ($+1$) & $24.01^\circ$ & $\pm 45^\circ$ & AAOS (1965) \\
\addlinespace[2pt]
Neck Ry \newline \textit{Universal} & Trunk extension ($+1$) & $7.67^\circ$ & $\pm 45^\circ$ & AAOS (1965) \\
\addlinespace[2pt]
Spine Rx \newline \textit{Universal} & Trunk lateral flexion ($+1$) & $0.00^\circ$ & $-25^\circ \dots +80^\circ$ & AAOS thoracolumbar \\
\addlinespace[2pt]
Spine Ry \newline \textit{Universal} & Trunk lateral bend ($+1$) & $0.00^\circ$ & $\pm 35^\circ$ & AAOS thoracolumbar \\
\addlinespace[2pt]
Torso Rz \newline \textit{Revolute} & Axial thorax twist ($+1$) & $-26.81^\circ$ & $\pm 60^\circ$ & Cheetham et al. (2008) \\
\addlinespace[2pt]
Lead Elbow \newline \textit{Revolute} & Elbow flexion ($+1$) & $0.00^\circ$ & $-5^\circ \dots +150^\circ$ & AAOS (1965) \\
\addlinespace[2pt]
Trail Elbow \newline \textit{Revolute} & Elbow flexion ($-1$) & $0.00^\circ$ & $-5^\circ \dots +150^\circ$ & AAOS (1965) \\
\addlinespace[2pt]
Forearms \newline \textit{Revolute} & Forearm pronation ($+1$ L / $-1$ R) & $90.00^\circ$ & $\pm 80^\circ$ & AAOS (1965) \\
\addlinespace[2pt]
Scapulae \newline \textit{Universal} & Protraction / Elev. ($+1$ L / $-1$ R) & $0.00^\circ$ & $\pm 25^\circ / -10^\circ \dots 40^\circ$ & Ludewig et al. (2009) \\
\addlinespace[2pt]
Knees Rz \newline \textit{Revolute} & Knee flexion ($-1$) & $0.00^\circ$ & $-5^\circ \dots +135^\circ$ & Winter (2009) \\
\addlinespace[2pt]
Ankles Rx/Ry \newline \textit{Universal} & Inversion / Dorsiflex. ($-1/+1, -1$) & $0.00^\circ$ & $-15^\circ \dots 35^\circ / -50^\circ \dots 20^\circ$ & AAOS (1965) \\
\addlinespace[2pt]
Midfoot Rz \newline \textit{Revolute} & Toe extension ($-1$) & $0.00^\circ$ & $-45^\circ \dots +70^\circ$ & AAOS first MTP \\
\bottomrule
\end{tabularx}
\end{table}

\section{Multibody Topology, Block Budget, and Variant Ladder}

\subsection{Compiled Block Budget and License Governance}
Execution of multibody golf swing models within MATLAB R2025b is bounded by the strict non-virtual block constraint of the Home license:
\begin{itemize}
    \item \textbf{Theoretical Hard Ceiling:} Exactly 1,000 compiled non-virtual blocks. Probed synthetically via \texttt{gs3dx\_license\_limit\_probe.m}, 1,000 blocks compile and simulate; 1,001 blocks abort with an unrecoverable Simulink compiler error.
    \item \textbf{Compiled Count Evaluation:} Evaluated exclusively on the compiled system:
    \begin{equation}
        N_{\text{compiled}} = \text{length}\left(\text{find\_system}(mdl, \text{'Compiled'}, \text{'on'}, \text{'Virtual'}, \text{'off'})\right).
    \end{equation}
    \item \textbf{Mandatory Validation Reserve:} All models must maintain a 25-block reserve, enforcing a hard upper ceiling of \textbf{975 compiled blocks} to accommodate transient in-memory instrumentation blocks added by diagnostic harnesses (e.g., \texttt{gs3dx\_contact\_check}, which appends 10 compiled sensing blocks).
\end{itemize}

\subsection{The GS3DX Model Progression Ladder}
The architectural lineage spans thirteen programmatically constructed variants. Entries record historical diagnostic evidence from the turnover handoff, not physical qualification:

\begin{table}[!htb]
\centering
\footnotesize
\caption{The Simscape Multibody GS3DX Structural Model Progression Ladder}
\label{tab:variant_ladder}
\begin{tabularx}{\textwidth}{p{4.2cm} >{\raggedright\arraybackslash}X c >{\raggedright\arraybackslash}X}
\toprule
\textbf{Model Variant / Builder} & \textbf{Structural Modifications} & \textbf{Blocks} & \textbf{Historical Diagnostic Evidence} \\
\midrule
\texttt{GS3DX\_Baseline} \newline \textit{\texttt{gs3dx\_clone\_baseline}} & Verbatim renamed clone of legacy model & 941 & Match on 413 signals across 344 steps \\
\addlinespace[2pt]
\texttt{GS3DX\_Slim} \newline \textit{\texttt{gs3dx\_build\_slim}} & Direct \texttt{InputTorque} drive; drops sources & 773 & Matches baseline to $6 \times 10^{-14}$~m clubhead diff \\
\addlinespace[2pt]
\texttt{GS3DX\_Quat} \newline \textit{\texttt{gs3dx\_build\_quat}} & Quaternion Spherical shoulders and 6-DOF hip & 740 & Eliminates joint lock; 7--9\% fewer steps \\
\addlinespace[2pt]
\texttt{GS3DX\_FullBody} \newline \textit{\texttt{gs3dx\_build\_lower\_body}} & Lower body legs added; feet welded to World & 945 & Passive legs; knee hyperextends ($+67^\circ$) \\
\addlinespace[2pt]
\texttt{GS3DX\_FullBodyContact} \newline \textit{\texttt{gs3dx\_build\_contact}} & 3 sole spheres/foot on plane; free pelvis & 967 & Stands from rest (slip $\le 2.1$~mm; lift 0~mm) \\
\addlinespace[2pt]
\texttt{GS3DX\_Golfer} \newline \textit{\texttt{gs3dx\_build\_golfer}} & de Leva 80~kg masses assigned to 15 solids & 967 & Sensed mass 80.393~kg; stands from rest \\
\addlinespace[2pt]
\texttt{GS3DX\_Fit} \newline \textit{\texttt{gs3dx\_build\_fit}} & Capture lengths; fitted hand-grip standoffs & 967 & Wrist IK errors reduced from 88 to 18~mm \\
\addlinespace[2pt]
\texttt{GS3DX\_FitLegs} \newline \textit{\texttt{gs3dx\_build\_fit\_legs}} & Leg servo reference trajectories from IK & 967 & Stance holds under passive upper-body drive \\
\addlinespace[2pt]
\texttt{GS3DX\_FitTrack} \newline \textit{\texttt{gs3dx\_build\_fit\_track}} & Feedforward + PD on 12 upper-body joints & 967 & Angle error $0.25^\circ$; pelvis tips without balance \\
\addlinespace[2pt]
\texttt{GS3DX\_FitBalance} \newline \textit{\texttt{gs3dx\_build\_fit\_balance}} & Damped inverse-Jacobian COM balance loop & 973 & Pelvis error at impact reduced 503 to 73~mm \\
\addlinespace[2pt]
\texttt{GS3DX\_Shape} \newline \textit{\texttt{gs3dx\_build\_shape}} & de Leva gyration inertia; visual ellipsoids & 973 & Vertical COM error reduced to 2.9~mm RMS \\
\addlinespace[2pt]
\texttt{GS3DX\_Neck} \newline \textit{\texttt{gs3dx\_build\_neck}} & 2-axis Universal neck driven by head markers & 975 & Vertical head error reduced 52.2 to 36.8~mm \\
\addlinespace[2pt]
\texttt{GS3DX\_Human} \newline \textit{\texttt{gs3dx\_build\_human}} & 5 contacts/foot; sprung midfoot; STL mesh & 965 & Square face at address; unstable finish dynamics \\
\bottomrule
\end{tabularx}
\end{table}

\subsection{Quaternion-to-Euler Mapping and Torque Duality}
\texttt{GS3DX\_Quat} replaces Gimbal joints with native Spherical Joints parameterized by unit quaternions $\mathbf{q} = [q_w, q_x, q_y, q_z]^T$. As documented by MathWorks~\cite{MathWorksSpherical}, the underlying joint primitive eliminates quaternion orientation singularities. However, to interface with legacy Euler-angle signal buses, \texttt{gs3dx\_xyz\_map.m} extracts intrinsic $X$-$Y$-$Z$ angles using the exact sequence:
\begin{equation}
    \mathbf{R} = \mathbf{R}_x(a) \mathbf{R}_y(b) \mathbf{R}_z(c).
\end{equation}
Let $\mathbf{R} \in \mathbb{R}^{3 \times 3}$ denote the rotation matrix computed from unit quaternion $\mathbf{q}$. The extracted angles are evaluated as:
\begin{align}
    a &= \text{atan2}(-R_{23}, R_{33}), \\
    b &= \text{atan2}(R_{13}, \text{hypot}(R_{11}, R_{12})), \\
    c &= \text{atan2}(-R_{12}, R_{11}),
\end{align}
with $a$ and $c$ unwrapped to the continuous branch of reference angles $\mathbf{q}_{\text{ref}}$.
\begin{quote}
\textbf{Singularity Caveat:} The adapter mapping remains \textbf{singular where $\cos(b) = 0$ ($|b| = 90^\circ$)}. While the Spherical Joint itself does not lock, representing motion through this adapter re-introduces the kinematic singularity; full singularity removal is not achieved.
\end{quote}

The kinematic velocity mapping relates Euler angle rates $\dot{\mathbf{q}}_r = [\dot{a}, \dot{b}, \dot{c}]^T$ (in radians/s) to the follower-frame angular velocity $\boldsymbol{\omega}_F$:
\begin{equation}
    \boldsymbol{\omega}_F = \mathbf{E}(b, c) \dot{\mathbf{q}}_r, \quad \text{where} \quad \mathbf{E}(b, c) = \begin{bmatrix} \cos b \cos c & \sin c & 0 \\ -\cos b \sin c & \cos c & 0 \\ \sin b & 0 & 1 \end{bmatrix}.
\end{equation}
By the principle of virtual power, mechanical power is invariant under generalized coordinates: $\mathbf{T}^T \boldsymbol{\omega}_F = \boldsymbol{\tau}_q^T \dot{\mathbf{q}}_r$. The torque dual transformation mapping Gimbal axis torques $\boldsymbol{\tau}$ and joint damping to follower-frame torque $\mathbf{T}$ is:
\begin{equation}
    \mathbf{T} = \mathbf{E}^{-T} \left(\boldsymbol{\tau} - \mathbf{D} \dot{\mathbf{q}}_{\text{deg}}\right),
\end{equation}
accounting explicitly for conversion between degrees and radians ($\dot{\mathbf{q}}_{\text{deg}} = \dot{\mathbf{q}}_r \times 180 / \pi$).

\section{Anthropometry, Mass Distribution, and Inertial Calibration}

\subsection{Standardized Segment Masses and Inertia Tensors}
The model's mass distribution is parameterized according to de Leva's (1996)~\cite{deLeva1996} anatomical adjustments to Zatsiorsky-Seluyanov's gamma-ray scanning standards for adult males. For an 80.0~kg golfer, the segments, masses, center of mass fractions, and central radii of gyration are detailed in Table~\ref{tab:de_leva_masses}. Total system mass equals $80.393$~kg, which incorporates 80.000~kg of human biological tissue plus 0.393~kg of equipment (0.330~kg club shaft and head, 0.057~kg grip assembly, and 0.006~kg contact sphere representations).

\begin{table}[!htb]
\centering
\small
\caption{Segment Anthropometry, Mass Allocation, and Inertial Parameters for 80~kg Golfer}
\label{tab:de_leva_masses}
\begin{tabularx}{\textwidth}{lcccccX}
\toprule
\textbf{Segment} & \textbf{Fraction} & \textbf{Mass (kg)} & \textbf{COM ($z/L$)} & \textbf{$r_{xx}/L$} & \textbf{$r_{yy}/L$} & \textbf{$r_{zz}/L$} \\
\midrule
Head + Neck & 0.0694 & 5.552 & 0.5002 & 0.303 & 0.315 & 0.261 \\
Thorax (Upper Trunk) & 0.1596 & 12.768 & 0.5066 & 0.505 & 0.449 & 0.347 \\
Abdomen (Mid Trunk) & 0.1633 & 13.064 & 0.3989 & 0.482 & 0.383 & 0.468 \\
Pelvis (Lower Trunk) & 0.1117 & 8.936 & 0.6014 & 0.615 & 0.551 & 0.587 \\
Thigh (each) & 0.1416 & 11.328 & 0.4095 & 0.329 & 0.329 & 0.149 \\
Shank (each) & 0.0433 & 3.464 & 0.4459 & 0.282 & 0.282 & 0.092 \\
Foot Rearfoot (each) & 0.0106 & 0.846 & 0.4415 & 0.257 & 0.245 & 0.124 \\
Foot Forefoot (each) & 0.0031 & 0.250 & 0.5000 & 0.200 & 0.200 & 0.100 \\
Upper Arm (each) & 0.0271 & 2.168 & 0.5772 & 0.285 & 0.269 & 0.158 \\
Forearm (each) & 0.0162 & 1.296 & 0.4574 & 0.276 & 0.265 & 0.121 \\
Hand (each) & 0.0061 & 0.488 & 0.7900 & 0.628 & 0.513 & 0.332 \\
\bottomrule
\end{tabularx}
\end{table}

\subsection{Parametric Segment Dimensions From Capture Geometry}
Using \texttt{gs3dx\_capture\_joint\_centres.m}, segment lengths are extracted from marker triad centroids and mapped into parameter variables ($Fit*$) that cannot be overwritten by drive files:
\begin{itemize}
    \item Thigh length ($L_{\text{thigh}}$): 0.460~m (hip joint center to knee axis).
    \item Shank length ($L_{\text{shank}}$): 0.424~m (knee axis to ankle joint center).
    \item Forearm length ($L_{\text{forearm}}$): 0.280~m (elbow joint center to wrist joint center).
    \item Upper arm length ($L_{\text{arm}}$): 0.305~m (shoulder joint center to elbow axis).
    \item Pelvis-to-shoulder trunk line ($L_{\text{trunk}}$): 0.490~m.
\end{itemize}

\subsection{Midfoot Joint Architecture and Mass Partitioning}
In \texttt{GS3DX\_Human}, a sprung revolute joint is introduced at the metatarsophalangeal (MTP) boundary, positioned at 73\% of foot length from the posterior heel and 25~mm above the sole:
\begin{equation}
    \tau_{\text{MTP}} = -K_{\text{midfoot}} \theta_{\text{MTP}} - C_{\text{midfoot}} \dot{\theta}_{\text{MTP}},
\end{equation}
where $K_{\text{midfoot}} = 2,000$~N$\cdot$m/rad and $C_{\text{midfoot}} = 0.5$~N$\cdot$m$\cdot$s/rad. Total foot mass (1.096~kg) is partitioned into a 0.250~kg forefoot solid carrying the toe contacts and a 0.846~kg rearfoot solid carrying the heel and ball contacts.

\section{Motion Capture Calibration, Inverse Kinematics, and Identifiability}

\subsection{Capture Pre-Processing and Trajectory Gaps}
The reference motion (\texttt{capture-A}) is sampled at 360~Hz over 654 frames ($T = 1.814$~s). Key temporal events are:
\begin{itemize}
    \item \textbf{Top of Backswing:} Frame 378 ($t = 1.050$~s), defined by pelvis yaw reversal.
    \item \textbf{Peak Clubhead Speed:} Frame 476 ($t = 1.319$~s), reaching 50.7~m/s approximately 16~cm before ball impact.
    \item \textbf{Ball Contact:} Frame 477 ($t = 1.322$~s), when the clubhead returns to address position along the target line; head speed drops to 41.4~m/s across contact.
\end{itemize}
Marker dropouts occur on the clubhead at frame 445, frames 519--551, and frames 653--654. The inverse kinematics algorithm omits gap-filled samples from the residual vector (\texttt{jc.gap}); fitting linearly interpolated points introduces severe tracking distortion ($> 0.5$~m) through follow-through.

\subsection{Levenberg--Marquardt Inverse Kinematics Formulation}
Inverse kinematics minimizes residual tracking errors over 33 independent coordinates $\mathbf{q} \in \mathbb{R}^{33}$ using \texttt{KinematicsSolver} as forward kinematics:
\begin{equation}
    \min_{\mathbf{q}_k} \frac{1}{2} \|\mathbf{r}_{\text{marker}}(\mathbf{q}_k)\|^2 + \frac{w_p^2}{2} \|\mathbf{q}_k - \mathbf{q}_{\text{posture}}\|^2 + \frac{w_s^2}{2} \|\mathbf{q}_k - \mathbf{q}_{k-1}\|^2 + \frac{w_{\text{rom}}^2}{2} \|\mathbf{r}_{\text{rom}}(\mathbf{q}_k)\|^2,
\end{equation}
where $w_p = 0.01$~m/rad pulls trunk coordinates toward neutral, $w_s = 0.02$~m/rad penalizes frame-to-frame coordinate difference (note: this is a coordinate step penalty, not velocity without division by $\Delta t$), and $w_{\text{rom}} = 3.0$~m/rad penalizes joint range violations:
\begin{equation}
    r_{\text{rom}, j}(\mathbf{q}) = \max\left(0, A_j(\mathbf{q}) - A_{j,\max}\right) + \max\left(0, A_{j,\min} - A_j(\mathbf{q})\right).
\end{equation}

\subsection{Joint Angle and Standoff Identifiability Limits}
\begin{enumerate}
    \item \textbf{Lead Elbow Flexion Identifiability:} Fitting unconstrained skin markers yields an unphysiological lead elbow flexion of $51.8^\circ$--$65.9^\circ$. When marker offsets are calibrated against an unconstrained fit, they absorb marker error; enforcing an anatomical ROM penalty thereafter causes severe conflict, inflating tracking RMS from 23.3~mm to 43.2~mm. Offsets must be calibrated under the constrained prior.
    \item \textbf{Hand-on-Grip Standoff Identifiability:} Total hand separation across the shaft ($72$--$74$~mm) is identifiable; however, the lateral shift of the shaft relative to the hands is unidentifiable because sideways shaft movement is absorbed by clubhead marker offsets. The standoff is partitioned symmetrically ($1.46$~in on each side).
\end{enumerate}

\section{Control Architecture: Feedforward, Feedback, and Balance}

\subsection{Upper-Body Feedforward and PD Tracking}
Upper-body actuation operates under \texttt{UpperBodyTracking} with an unactuated 6-DOF pelvis (\texttt{ModelingMode} 0):
\begin{equation}
    \boldsymbol{\tau}_{\text{upper}}(t) = \boldsymbol{\tau}_{\text{ff}}(t) + \mathbf{K}_p (\mathbf{q}_{\text{ref}}(t) - \mathbf{q}(t)) + \mathbf{K}_d (\dot{\mathbf{q}}_{\text{ref}}(t) - \dot{\mathbf{q}}(t)).
\end{equation}
In \texttt{gs3dx\_track\_gains.m}, gains are critically damped ($\zeta = 1.0$) at natural frequency $\omega_n = 2\pi \times 6.0$~rad/s based on reflected segment inertia $I_j$ (in kg$\cdot$m$^2$). The continuous radian gains are converted to Simscape degree units by multiplying by $\pi / 180$:
\begin{equation}
    K_{p, j} = I_j \omega_n^2 \left(\frac{\pi}{180}\right), \quad K_{d, j} = 2 \zeta I_j \omega_n \left(\frac{\pi}{180}\right).
\end{equation}

\subsection{Iterative Learning Control (ILC) Formulation}
Feedforward torques $\mathbf{F}(t)$ are updated via a filtered P-D learning law:
\begin{equation}
    \mathbf{F}^{(k+1)}(t) = \mathcal{Q}\left(\mathbf{F}^{(k)}(t) + \gamma \cdot \left(\mathbf{K}_p \mathbf{e}^{(k)}(t) + \mathbf{K}_d \dot{\mathbf{e}}^{(k)}(t)\right)\right),
\end{equation}
where $\mathcal{Q}$ is a zero-phase 4th-order Butterworth low-pass filter (cutoff 12~Hz), $\gamma$ is the learning gain (default 1.0), $\mathbf{e}^{(k)} = \mathbf{q}_{\text{ref}} - \mathbf{q}^{(k)}$, and $\dot{\mathbf{e}}^{(k)} = \dot{\mathbf{q}}_{\text{ref}} - \dot{\mathbf{q}}^{(k)}$.

\begin{quote}
\textbf{Convergence Caveat:} Because $\mathcal{Q}$ attenuates high-frequency content, an ILC fixed point $\mathbf{F}^* = \mathcal{Q}(\mathbf{F}^* + \gamma \mathbf{u}_{\text{pd}}^*)$ does \textbf{not} imply that PD torque vanishes. In historical tests, Iteration 2 exhibited the lowest observed PD torque ($10.3$~N$\cdot$m RMS), but is not a proven global optimum. Later iterations exhibited growing PD torque (14.7~N$\cdot$m at iteration 3; 20.7~N$\cdot$m at iteration 4).
\end{quote}

\subsection{Leg Servo and Damped Inverse-Jacobian Balance Feedback}
The lower body is actuated by twelve joint coordinates $\mathbf{q}_{\text{leg}} \in \mathbb{R}^{12}$. The balance gain $\mathbf{G}(t)$ is computed in \texttt{gs3dx\_balance\_gain.m} from the per-foot $6 \times 6$ foot pose Jacobian $\mathbf{J}$ (position in m, rotation vector in rad scaled by $\text{lever\_m} = 0.5$~m), using central finite differences in degrees ($h = 10^{-3}$ deg). Damped least-squares inversion with damping $\lambda = 0.002$~m/deg yields:
\begin{equation}
    \mathbf{G}_{\text{leg}}(t) = -\mathbf{J}^T \left(\mathbf{J} \mathbf{J}^T + \lambda^2 \mathbf{I}\right)^{-1} \begin{bmatrix} \mathbf{I}_D \\ \mathbf{0}_{6-D, D} \end{bmatrix},
\end{equation}
where $D = 3$ (or 2 for horizontal axes). In \texttt{gs3dx\_balance\_command.m}, the commanded torque vector is:
\begin{equation}
    \mathbf{cmd} = \mathbf{C}_0 + \mathbf{K}_p \odot \left(\mathbf{A}(t) + \mathbf{G}(t) \mathbf{s} + \mathbf{foot}\right) + \mathbf{K}_d \odot \left(\mathbf{R}(t) + \mathbf{G}(t) \mathbf{r}_s\right),
\end{equation}
with shift $\mathbf{s} = \text{sat}_{\text{limit}}(-k_p \mathbf{e} - k_d \dot{\mathbf{e}})$, separately commanded rate $\mathbf{r}_s = -k_p \dot{\mathbf{e}}$, and separate foot displacement. This rate is not the exact derivative of the saturated shift: it omits acceleration and saturation derivatives. The implementation uses this approximation:
\begin{equation}
    \mathbf{foot}_{\text{leg}} = \mathbf{G}_{\text{leg}}(t) \cdot \text{sat}_{\text{limit}}\left(k_f \mathbf{e}_{\text{foot}}\right).
\end{equation}
Builder defaults are $k_p = 3.0$, $k_d = 0.4$, $k_f = 1.0$, and $\text{limit} = 0.10$~m.

\section{Theoretical Caveats: Nullspaces and Under-Actuation}

\subsection{Closed-Loop Internal Force Nullspace Hypothesis}
Let $\mathbf{J}_c \in \mathbb{R}^{6 \times n}$ denote the constraint Jacobian governing the grip closure. Internal constraint torques satisfy:
\begin{equation}
    \boldsymbol{\tau}_{\text{int}} = \mathbf{J}_c^T \boldsymbol{\lambda} \in \text{range}(\mathbf{J}_c^T).
\end{equation}
Let $\mathbf{Z}$ denote a matrix whose columns span $\text{null}(\mathbf{J}_c)$ (so $\mathbf{J}_c \mathbf{Z} = \mathbf{0}$). For any constraint-compatible generalized velocity $\mathbf{v} = \mathbf{Z} \boldsymbol{\nu}$:
\begin{equation}
    \mathbf{Z}^T \boldsymbol{\tau}_{\text{int}} = \mathbf{Z}^T \mathbf{J}_c^T \boldsymbol{\lambda} = (\mathbf{J}_c \mathbf{Z})^T \boldsymbol{\lambda} = \mathbf{0}.
\end{equation}
For quaternion configurations, $\dot{\mathbf{q}}=\mathbf{N}(\mathbf{q})\mathbf{v}$; configuration and velocity dimensions may differ. The constraint Jacobian acts on velocity. Full multibody dynamics include actuation $\mathbf{B} \mathbf{u}$ and ground contact wrenches:
\begin{equation}
    \mathbf{M}(\mathbf{q}) \dot{\mathbf{v}} + \mathbf{h}(\mathbf{q}, \mathbf{v}) = \mathbf{B} \mathbf{u} + \mathbf{J}_c^T \boldsymbol{\lambda} + \mathbf{J}_{\text{contact}}^T \mathbf{w}_{\text{contact}}.
\end{equation}
To generate a selected internal torque using commanded actuators alone, that
torque must lie in the actuated image $\text{range}(\mathbf{B})$. Passive
constraint reactions are not independent control inputs.
\begin{quote}
\textbf{Theoretical Note:} The explanation of torque drift as accumulation within closed-loop kinematic nullspaces represents a \textbf{plausible scientific hypothesis}, supported by simulation observations, but is \textbf{not an established mathematical proof}.
\end{quote}

\subsection{Rank-Deficiency vs. Global Dynamic Infeasibility}
\begin{quote}
\textbf{Mathematical Caveat:} The instantaneous \textbf{rank-deficiency of contact wrench or kinematic Jacobian matrices does not prove the global infeasibility} of the forward-dynamics swing matching problem.
\end{quote}
Local Jacobian rank-deficiency demonstrates the limitation of static linear inversion schemes, but does not establish global dynamic infeasibility. Conversely, unproven claims of Lie algebraic nonholonomic controllability must be avoided without rigorous proof.

\section{Contact Wrench Mechanics, Support Hull, and Ground Reactions}

\subsection{Planar Support Hull and Point-Force Contact Diagnostic}
The point-force diagnostic in \texttt{gs3dx\_contact\_wrench.m} projects demanded center of mass accelerations onto the planar foot support hull under assumed Coulomb friction $\mu = 0.8$. Dynamic equilibrium equations are:
\begin{align}
    \sum_{i} \mathbf{f}_i + M \mathbf{g} &= M \mathbf{a}_{\text{COM}}, \\
    \sum_{i} (\mathbf{p}_i - \mathbf{r}_{\text{COM}}) \times \mathbf{f}_i + \sum_{i} \boldsymbol{\tau}_{\text{contact}, i} &= \dot{\mathbf{h}}_{\text{COM}}.
\end{align}
\begin{quote}
\textbf{Diagnostic Scope:} The frozen residuals (\textbf{21.863~N$\cdot$m RMS through impact; 81.427~N$\cdot$m at impact}) represent a \textbf{point-force diagnostic on a planar support polygon}, omitting ground reaction yaw couples and full 3D segment angular momentum rates. Passing this diagnostic does \textbf{not} establish complete dynamic admissibility.
\end{quote}
Contact normal force formulations in Simscape utilize internal smoothing; formulas $F_n = \max(0, -K_c z - C_c \dot{z})$ are schematic linear approximations.

\subsection{Measured Ground Reactions vs. Kinematic GRF Estimates}
\begin{itemize}
    \item \textbf{No Empirical Force Plates:} The capture source file contains zero analog force plate channels.
    \item \textbf{Kinematic GRF Estimates:} Total GRF is estimated from segment accelerations:
    \begin{equation}
        \mathbf{F}_{\text{kin\_GRF}}(t) = \sum_{i} m_i (\ddot{\mathbf{r}}_i(t) - \mathbf{g}) + m_{\text{club}} (\ddot{\mathbf{r}}_{\text{club}}(t) - \mathbf{g}).
    \end{equation}
    At address, $\mathbf{F}_{\text{kin\_GRF}} = 1.001$~BW. Prior to impact, filtered peak is $1.24$--$1.33$~BW (raw peak $1.853$~BW at $-55.6$~ms). Club acceleration contributes $+0.57$~BW. Contact compliance parameters remain uncalibrated assumptions.
\end{itemize}

\section{Stability Analysis: Distinguishing Baselines From Unstable Dynamics}

\subsection{Stable Baselines vs. Current Unstable Human Dynamics}
\begin{itemize}
    \item \textbf{Stable Baselines:} \texttt{GS3DX\_Baseline}, \texttt{GS3DX\_Slim}, and \texttt{GS3DX\_Quat} integrate the impact drive cleanly across 344 steps; \texttt{GS3DX\_FullBodyContact} and \texttt{GS3DX\_Golfer} stand stably from rest (slip $\le 2.1$~mm, lift 0~mm, Newton residual 0.71~N$\cdot$s).
    \item \textbf{Current Unstable Dynamics:} Full-swing simulations of \texttt{GS3DX\_Human} exhibit transient instability through impact and finish.
\end{itemize}

\subsection{Rejection of Extreme Transient Force Spikes}
Historical simulation logs recorded extreme vertical contact force spikes:
\begin{enumerate}
    \item \textbf{2,466~BW Spike ($t = 0.386$~s):} Integrator stall at $h_{\min}$ due to derivative rate steps.
    \item \textbf{365~BW Spike ($t = 0.449$~s):}
    Simulation with active balance and leg servo produced a 365~BW ($287{,}500$~N) spike on the right lesser toe. At $t = 0.35$--$0.40$~s, the golfer was airborne. On row 12, servo ($-3{,}755.6$~N$\cdot$m) and balance ($+2{,}130.0$~N$\cdot$m) fought, driving the ankle into $42.5^\circ$ plantarflexion before singular toe impact. Near impact, support dropped to 0.0~BW (kinematic demand is 1.85~BW).
    \item \textbf{22.6~BW Spike ($t = 0.418$~s):}
    With row 12 balance gain zeroed, the right ankle dropped $0.661$~m in $0.060$~s before strike; midfoot angle stepped from $+0.013$ to $-5.452$~rad in 10~ms. Pelvis tilt diverged from $22.5^\circ$ to $101.9^\circ$. Refining maximum step size to $0.0001$~s yielded an identical 23.1~BW spike.
\end{enumerate}
Both the \textbf{365~BW and 22.6~BW readings reject physical qualification}. Controller conflict and airborne strike are suspected causes; the causal explanation is not yet experimentally established. The model does \textbf{not} possess an open-loop stable swing.

\section{Current Acceptance Gates, Open Work, and Red Thresholds}

\subsection{Gate 1: Lead Arm Straight Extension (\#11156)}
\begin{itemize}
    \item \textbf{Requirement:} Lead elbow flexion $\le 20.0^\circ$; tracking penalty $\le \text{free} + 15.0$~mm (ceiling $38.3$~mm).
    \item \textbf{Current Status: RED.} Free calibration yields $18.5^\circ$--$41.7^\circ$ flexion and 43.2~mm RMS; banded calibration yields $65.8^\circ$ flexion and 76.8~mm RMS. Thresholds cannot be loosened without owner sign-off.
\end{itemize}

\subsection{Gate 2: Clubhead Spatial-Temporal Matching (\#11160)}
\begin{itemize}
    \item \textbf{Requirement:} Defined in \texttt{gs3dx\_club\_match.m}:
    \begin{enumerate}
        \item Clubhead path: $\le 10.0$~mm RMS and $\le 20.0$~mm maximum.
        \item Shaft orientation: $\le 2.0^\circ$ maximum.
        \item Impact clubhead speed error: $\le 2.0\%$.
        \item Clubface orientation error at contact: $\le 2.0^\circ$.
    \end{enumerate}
    \item \textbf{Current Status: RED.} Actual values: head RMS 20.012~mm, head max 41.987~mm, shaft max $3.11^\circ$, speed error $-2.009\%$, wrist max 53.5 / 68.6~mm (limit 35~mm). Missing \texttt{phases.contact} keeps \texttt{.face} false and rejects the match; peak speed is not contact. 16 synthetic R2025b tests passed after 7 red.
\end{itemize}

\subsection{Gate 3: Dynamically Consistent Leg Reference (\#11173 Step 4)}
\begin{itemize}
    \item \textbf{Requirement:} Pelvis residual wrench $< 5\%$ body weight; impact support within $0.20$~BW of kinematic GRF.
    \item \textbf{Current Status: RED.} Support near impact is 0.0~BW. Leg servo ($327.6$~N$\cdot$m RMS) and balance ($314.3$~N$\cdot$m RMS) destructively cancel ($\approx 101$~N$\cdot$m net).
\end{itemize}

\section{Multi-Engine Evaluation and Subject Matching Workflows}

Under epic \#11161, biomechanics are evaluated across Simscape, MuJoCo, Drake, Pinocchio, and OpenSim:
\begin{itemize}
    \item \textbf{MuJoCo Benchmark (\#11166):} Scaling \texttt{capture-A} to subject geometry (\texttt{capture-O}, 1.956~m, 104.3~kg) tracked with $0.060$~m drift, whereas raw \texttt{capture-O} diverged ($0.567$~m drift).
    \item \textbf{ZMP Support Margin:} This experiment indicated that for \texttt{capture-O}, the reference ZMP fell outside the physical support hull in 86\% of frames (vs. 47\% for \texttt{capture-A}). This is an attributed multi-engine observation; \texttt{capture-O} is not yet Simscape qualified.
\end{itemize}

\section{Visualization, Offscreen Rendering, and Asset Provenance}

\subsection{Headless Offscreen Graphics Architecture}
\texttt{gs3dx\_render.m} extracts Cartesian poses via \texttt{KinematicsSolver} and renders offscreen figures under batch execution:
\begin{itemize}
    \item \textbf{Face-On View:} Camera on $+X_W$ (azimuth $90^\circ$, elevation $5^\circ$).
    \item \textbf{Down-the-Line View:} Camera on $-Y_W$ (azimuth $0^\circ$, elevation $5^\circ$).
    \item \textbf{Top View:} Camera on $+Z_W$ (azimuth $0^\circ$, elevation $90^\circ$).
\end{itemize}

\subsection{Parametric Driver Head Provenance}
Rendered from parametric binary STL \texttt{models/gs3dx\_driver\_head.stl} ($124 \times 115 \times 61$~mm, $10.5^\circ$ loft) attached as a massless File Solid with \texttt{FaceSquareRoll} $-20.73^\circ$.

\section{Operational Qualification Protocol and New-Agent Checklist}

\subsection{Process Concurrency and Safety Rules}
\begin{enumerate}
    \item \textbf{Serialized Runner Protocol:} Never launch MATLAB directly when the user home lock file (\texttt{gs3dx\_matlab.lock}) contains an active process ID. Always use the serialized runner script \texttt{exploratory\_gs3dx/tools/run\_matlab\_locked.ps1}.
    \item \textbf{Simulink Cache Redirection:} \texttt{gs3dx\_setup.m} redirects cache to \texttt{tempdir}. Never call \texttt{setup\_matlab\_environment.m} (\texttt{savepath} contamination).
    \item \textbf{Immutable Source Models:} Never load or write original models in \nolinkurl{../src/model/GolfSwing3D_Kinetic.slx}. All writes must be routed through \texttt{gs3dx\_save\_model.m}, enforcing compiled blocks $\le 975$.
\end{enumerate}

\subsection{New-Agent Execution Checklist}
\begin{enumerate}
    \item Verify working directory is within the isolated worktree checkout; pin current appropriate source commit hash.
    \item Verify MATLAB version is explicitly R2025b (\nolinkurl{C:/Program Files/MATLAB/R2025b/bin/matlab.exe}).
    \item Run unit safety tests via the serialized runner.
    \item Execute document contract tests:
    \begin{quote}\footnotesize\ttfamily
    python3 -m pytest \\
    docs/research/simscape\_matching\_reference/test\_simscape\_matching\_reference.py
    \end{quote}
    \item Check red gate blockers (\#11156, \#11160, \#11173) and do not attempt to loosen frozen numerical tolerances.
\end{enumerate}



\section{Reproducible Capture-Specific Export}

The entry is \texttt{gs3dx\_export\_batch.m}, invoked through
\texttt{run\_matlab\_locked.ps1}. Set
\nolinkurl{GS3DX_CAPTURE_ID} (neutral alias),
\nolinkurl{GS3DX_OUTPUT_DIR} (explicit destination),
\nolinkurl{CAPTURE_REGISTRY_REPO} (provider checkout),
\nolinkurl{MATLAB_PYTHON_EXE} (compatible Python with ezc3d), and, when needed,
\nolinkurl{CAPTURE_DATA_DIR} (private capture root).
Private data and checkpoints must not enter git.
\begin{verbatim}
# Run from the exploratory GS3DX directory.
$env:GS3DX_CAPTURE_ID = 'capture-O'
# Set output/provider/Python variables before this call.
$exportArgs = @('-Script', 'tools/gs3dx_export_batch.m',
          '-Log', "$env:TEMP/gs3dx-export.log",
          '-TimeoutMin', '60')
powershell -NoProfile -File tools/run_matlab_locked.ps1 @exportArgs
\end{verbatim}

This entry exports an inverse-kinematics visualization, not a forward dynamic
rollout. Both captures use segment lengths measured by
\texttt{gs3dx\_fit\_lengths}. The unsaved model workspace used for solving and
rendering is identical. Kinematic lengths do not identify body mass or inertia;
the owner capture's known mass is not established by this geometry fit.

For sample rate $f_s$, integer stride $s$, and capture index $k_i$:
\[
  t_i=(k_i-1)/f_s,\qquad f_{\mathrm{video}}=f_s/s.
\]
Indices must have constant spacing. A 360-Hz capture sampled every 12 frames
and a 240-Hz capture sampled every 8 frames both play at 30 fps.
Do not append a nonuniform final frame. The final sampled timestamp and full
capture endpoint are recorded separately. A video with $N$ frames has container
duration $N/f_{\mathrm{video}}$; the first-to-last sampled timestamp span is
$(N-1)/f_{\mathrm{video}}$.

The marker overlay is $3\times n_{\mathrm{targets}}\times N$ in the address
waist frame, subset in the same order as pose columns. Still requests refer to
capture frames and are mapped to pose columns; filenames record the actual
selected index. Peak speed is a distinct event from contact.

The private IK checkpoint includes capture, model and source hashes, fitted
geometry, and requested frames. A differing identity forces a fresh solve.
Source mutation during a solve rejects the checkpoint identity. A successful
checkpoint or render does not pass anatomical, club, contact, or dynamics gates.

Before sharing, verify the natural successful-exit receipt, capture/model/code
identity, qualification label, complete decode, codec, frame rate, and
representative address, backswing, contact-near, and final sampled frames.
Historical clips without recovered identity or physical timing retain that
uncertainty.

\subsection{Fresh Baseline Fits and Required Refinement}
Capture A is the tour reference; capture O is the owner's optical capture.
The October 1 fresh exports use capture-specific geometry in
\texttt{GS3DX\_Fit} and inverse-kinematics poses. Forward dynamics does not
drive these videos. Their mean and maximum sampled-frame measured-target RMS
errors are respectively 12.745748/40.208544~mm for A and
26.558095/116.576104~mm for O. These aggregate joint-centre and club-target
residuals are not the separate club-head acceptance metrics.

Both final native export jobs exited naturally with code zero. Complete H.264
decoding verified 55 tour frames and 46 owner frames at 30~fps. Native success
and successful decoding do not qualify the physical matching gates. The
scratch mirror reports an unknown Git commit; exact capture, model, exporter,
IK and sampling file hashes identify the executed inputs instead.

Future matches and shareable exports must use the refined ellipsoid reference
style. The current Fit graphics are retained as comparison baselines. The
backwards-looking tour feed/feet require checks of time ordering, camera,
handedness and foot frames. The owner's larger residuals require per-target
and per-phase diagnostics of marker definitions, calibrated offsets,
capture-specific lengths, spine/neck freedom and solver continuity.
Avatar appearance does not establish player skill. Migrating a model's graphics
must preserve or explicitly requalify frames, joints, masses and inertias;
kinematic equivalence cannot establish dynamic equivalence.


\section{Refined Human Matching: Orientation and Identifiability}
The former Fit pose solver assumed 33 independent coordinates; the Human
model has additional neck coordinates. Joint roles must therefore be resolved
from unique native block paths and validated joint schemas. A graphics
replacement cannot substitute for a compatible kinematic model. The refined
exporter selects Human explicitly, fits its longitudinal ellipsoid dimensions
to the capture, and records the exact solve and rendering dependency hashes.
Transverse artistic proportions and uncalibrated subject inertias remain
limitations. Native model tests are required in addition to pure helper tests.

The 14 position targets do not determine every segment orientation. In
particular, ankle position alone leaves foot rotation ambiguous; a forward
axis alone also leaves roll ambiguous. Raw ankle-to-toe direction includes
the height difference between markers and is not the shoe's native forward
axis. Let $a_f$ denote ankle position and $t_{i,f},t_{o,f}$ the toe markers.
Construct the proper marker triad
\[
 x_f=\frac{(t_{i,f}+t_{o,f})/2-a_f}{\|(t_{i,f}+t_{o,f})/2-a_f\|},\quad
 z_f=\frac{x_f\times(t_{i,f}-t_{o,f})}{\|x_f\times(t_{i,f}-t_{o,f})\|},\quad
 F_f=[x_f,\ z_f\times x_f,\ z_f].
\]
For explicitly measured address frames, obtain the proper polar mean
$F_0=U\operatorname{diag}(1,1,\det(UV^T))V^T$ from the SVD of their triad sum.
The target shoe orientation is
\[
 R_{t,f}=F_fF_0^TR_z(\psi_0).
\]
Here horizontal address yaw $\psi_0$ comes from the capture. A horizontal sole
at address is an explicit calibration assumption, not force or contact
evidence. Missing address calibration, degenerate triads and unmeasured
markers must fail closed or remain explicitly masked; interpolation must not
be described as measured orientation.

For each foot, append the nine-component residual
\[
 r_{R,f}=w_R\operatorname{vec}(R_{p,f}-R_{t,f})/\sqrt{2},\qquad
 \|r_{R,f}\|=2w_R\sin(\theta_f/2).
\]
The two feet contribute 18 residual components. The weight has units of
metres per normalized orientation error, so it is a modeling tradeoff with
the position residuals; it is not an anatomical accuracy threshold.
Backward yaw and an inverted sole both receive a nonzero penalty.
The reported 14-target position RMS excludes this orientation term, and
angular errors must be reported separately. Full foot orientation does not
resolve missing thorax, pelvis or head orientation observations.

An owner reverse-pass candidate illustrates why appearance alone cannot
select a match: its mean sampled frame RMS increased from 26.558 to
32.185~mm while its maximum decreased from 116.576 to 46.052~mm.
Its native execution exited naturally with code zero, but it was rejected
as a replacement because it did not improve both metrics. Separate
target diagnostics locate a sampled owner club-head error of 375.358~mm
at capture frame 89 and a tour trail-elbow error of 96.517~mm at frame 517.
These are fitting defects to investigate, not evidence of golfer skill.

Fresh Human calibration on frames 1, 13 and 25 gives position RMS
1.188, 1.195 and 1.208~mm with both foot angular errors below $0.7^\circ$.
The native probe exited naturally with code zero at
\texttt{2026-10-02T05:01:00.641030Z}. These are the same frames used for
calibration, not held-out evidence. In contrast, applying the orientation
term with inherited Fit offsets gave 51.263, 43.243 and 39.399~mm position
RMS and was rejected. The full refined exports must test generalization
beyond this address interval and report marker-offset magnitudes explicitly.
Native joint-role extraction reports 37 independent Human fit coordinates
with grip closure. The earlier erroneous count of 39 resulted from selecting
closed-loop variables by the old model's joint identifiers.

The original worker's damping-times-ten stance experiment also failed:
initial heel force rose to 15,513~N, support vanished by 0.60~s, and at 1~s
pelvis tilt was $117^\circ$ with pelvis height $-1.83$~m. A composed pelvis
up-axis tilt of $22.5^\circ$ is not an individual hip start angle; zeroing
one start angle is not a valid leveling operation. The original worker queued
a composed-pose check stopping at 0.02~s; no outcome or full-swing qualification
is asserted here. Preserve that work rather than duplicating its queue.

\subsection{Full-Swing Refined Exports and Comparison}
The final Human fits use the same measured captures, sampled frames and
capture-specific lengths as their Fit baselines. Natural successful native
solve receipts are dated \texttt{2026-10-02T05:15:54Z} for A and
\texttt{2026-10-02T05:23:22Z} for O. These times identify new execution,
not approval of the physical gates. The model SHA-256 begins
\texttt{9a26ee80761f5108}; the accompanying provenance records its full hash,
recorded solve dependencies, geometry and capture hashes. The scratch source
mirror has no verified Git revision; its reported commit remains unknown.

\begin{table}[!htb]
\centering
\caption{Measured-target position RMS in mm, over identical sampled frames.
Each maximum is a maximum frame RMS, not a maximum individual target error.}
\begin{tabular}{lrrrr}
\toprule
Capture & Fit mean & Human mean & Fit maximum & Human maximum \\
\midrule
A: tour & 12.746 & 13.188 & 40.209 & 40.121 \\
O: owner & 26.558 & 16.767 & 116.576 & 37.626 \\
\bottomrule
\end{tabular}
\end{table}
The owner match improves both position metrics. The tour match exchanges a
small increase in mean position error for calibrated foot orientations;
it is not an improvement of every metric. Left/right measured foot angular
errors have maxima $12.83^\circ/19.82^\circ$ for A and
$19.55^\circ/20.38^\circ$ for O. A supplies 55 measured foot samples per side;
O supplies 45 of 46, with the missing sample explicitly masked. These are
orientation-fit diagnostics, not new acceptance thresholds.

All four shareable clips use the Human ellipsoid model, H.264 encoding,
800 by 600 pixels and 30 fps. Complete decoding verifies 55 frames for A
and 46 for O; their first-to-last sampled spans are 1.80 and 1.50 seconds.
Final captions distinguish pose sample number from measured capture frame.
Two views per capture are supplied: face-on and down-the-line. The original
Fit clips remain available for comparison. The shareable package excludes
raw optical recordings, pose caches and local private paths.

The Human owner fit still has wrist-offset norms of approximately 55 and
62 mm and a club-target offset norm of 106 mm. These magnitudes require
review of grip geometry, target definitions and offset identifiability.
The decorative ground plane is not a measured ball-contact plane.
Head and neck orientation are not constrained by the present IK objective,
although the capture may contain relevant head markers. Neither anatomical
orientation accuracy nor subject-specific inertia follows from low position
RMS. Forward dynamics does not drive these videos; all physical matching
gates remain unqualified.
\subsection{Subject Physics and a Rejected Grip-Calibration Candidate}
Capture-fitted lengths alone do not update the subject's physical model.
The in-memory Human adapter accepts a declared positive subject mass and
whitelisted segment lengths, verifies native solid expressions and SI units,
and preserves joint definitions and parameter metadata. It never saves the
model. The shared de Leva catalog scales segment mass as $m_s=f_sM$ and
principal moments as $I_i=m_s(k_iL_s)^2$, with the documented axisymmetric
transverse averaging. Forearm halves use the parallel-axis construction
already specified above. Split-foot rear mass is total foot mass minus the
fixed structural forefoot mass; the rear centre offset compensates the
forefoot lever arm. This is a model construction assumption, not a measured
individual inertia tensor.

Eighteen native adapter/catalog tests pass with zero incomplete tests and a
natural successful exit at \texttt{2026-10-02T06:03:51Z}. They cover baseline
parameters, subject mass scaling, length scaling, moment triangle inequalities,
forearm lumping, parameter metadata, unsupported fields and units, and
unchanged joint expressions. Equipment mass is obtained through the existing
solid inertia audit, which respects native inertia mode. A Custom zero-mass
visual must not be classified using an inactive density parameter. These
tests establish parameter consistency; they do not establish contact
equilibrium or a qualified forward-dynamics swing.

Functional grip calibration treats the wrist marker's trajectory in the
club-marker frame as a sphere. After centering and scaling the samples,
solve the algebraic least-squares system
\[
 [2P^T\;\mathbf{1}]
 \begin{bmatrix}c\\w\end{bmatrix}
 = \bigl[\|p_1\|^2,\ldots,\|p_n\|^2\bigr]^T,
 \qquad r^2=w+\|c\|^2.
\]
The design matrix must have numerical rank four; coplanar arcs, repeated
points and nonfinite input fail closed. Recover physical units after solving
and report radial residuals and matrix conditioning. The algebraic fit is
not a guarantee of unbiased anatomical joint-centre recovery. Fifteen pure
sphere tests and twelve controlled grip-contract tests pass with no incomplete
tests. Grip and marker-offset calibration samples are recorded separately;
reference club frames and shaft-axis estimates use only calibration samples.
Missing markers remain excluded. Equal bilateral transverse standoff is an
explicit gauge assumption.

The owner's functional-grip candidate used frames 1, 9,\ldots,169 for grip
calibration and retained 24 later sampled frames for evaluation. Native
execution exited naturally with code zero at \texttt{2026-10-02T06:15:53Z}.
Held-out mean/maximum frame RMS worsened from 26.011/37.626 to
27.969/41.278~mm; full-swing values worsened from 16.767/37.626 to
17.948/41.278~mm. The candidate was rejected and did not replace the best
Human clips, even though its lead-wrist offset norm dropped to 15.6~mm.
Capture lengths were shared estimates from the complete capture; held-out
status here applies to grip and offset calibration, not anthropometry.
The candidate applied the declared 104.3-kg subject mass and evaluated
0.3929236~kg of equipment, but mass does not drive the IK objective. No
improved dynamics or playing ability is inferred from this experiment.
\subsection{Fresh Contact Geometry and Imported Event Semantics}
A read-only native forward-kinematics audit of the selected owner address
pose completed with natural exit zero at \texttt{2026-10-02T06:42:14Z}.
It used fitted Human lengths and the declared 104.3-kg subject mass.
Against the existing model plane, the ten contact-sphere clearances ranged
from $-64.399$ to $-48.619$~mm, a 15.780-mm spread. Negative clearance
denotes geometric penetration, not a force measurement. No dynamics was
run, no simulation initial targets were changed, and no model was saved.
Ground placement and a consistent contact/gravity initialization therefore
remain separate verification steps. The pelvis frame's local up vector
tilted 57.889 degrees in this pose; this frame-relative quantity does not
establish an anatomical torso angle or a cause of instability.

Both imported recordings declare \texttt{POINT:UNITS} as metres and have no
\texttt{EVENT} parameter group. Under ezc3d 1.7.2, capture A has 717
negative-residual samples, all already nonfinite in XYZ; capture O has no
negative residuals but 1,173 nonfinite XYZ samples. Counts cover all marker
tracks. The fourth point row is homogeneous XYZ1; invalidity belongs to
\texttt{data.meta\_points.residuals}, not that row\cite{ezc3d}.
Neither peak club-head speed nor an address-line crossing establishes a
measured ball-contact event. Thirteen importer-contract tests, seventeen
runtime-fingerprint tests and twelve exporter-contract tests pass with zero
incomplete tests. Live component versions are MATLAB R2025b Update 5,
Simulink/Simscape/Multibody 25.2, Python 3.12.10, NumPy 2.4.4 and ezc3d 1.7.2.
The importer validates declared units, converts supported metres/centimetres/
millimetres to SI, and masks an entire marker sample when XYZ or residual
metadata is invalid. Required recorded source hashes fail closed; direct
runtime components are recorded without claiming all transitive packages or
external STL assets are frozen. Fresh exports exited naturally with code zero
at \texttt{2026-10-02T07:04:47Z} (A) and \texttt{07:13:40Z} (O).
All four MP4s fully decoded, with 55 tour and 46 owner samples; cache joint
arrays match the earlier selected Human fits within $10^{-12}$, with identical
capture/model hashes, geometry and frames. A new shareable ZIP contains these
clips and sanitized fresh runtime/source provenance. The earlier packages
remain available. This is an evidence update, not a new numerical improvement.

\subsection{Integration With Current Repository Models}
The continuation was committed as \texttt{69b4bd99f} after normal repository
hooks passed. Current main at \texttt{cee0a65e0} includes subsequent compiled
home-budget diagnostics, model promotion and anatomical meshes. Those changes
are preserved during integration. Its Human model SHA256 begins
\texttt{91997471}; the verified Desktop matches use the separately identified
\texttt{9a26ee80} model. Their successful receipts do not qualify the newer
model's full-swing fit. The integrated source passed 187 native MATLAB tests
with zero failed or incomplete tests, including Human visual/physical parameter
checks, and 44 focused Python checks. Full-swing matching against that model
remains required before replacing the selected outputs. The historical-player capture
handoff is preserved as a separate continuation scope in the canonical record.

\subsection{Current Model Comparison and Publication}
The reviewed continuation is published as draft PR
\href{https://github.com/D-sorganization/UpstreamDrift/pull/11256}{\#11256}.
Integration commit \texttt{7a315dd32} preserves current main through
\texttt{80dd2f7be}. The tour export using Human model \texttt{91997471}
exited naturally with code zero at \texttt{2026-10-02T07:49:26Z}; its
mean/max frame position RMS is 13.366/41.385 mm. Both are worse than the
selected \texttt{9a26ee80} Human result (13.188/40.121 mm), so this candidate
does not replace the selected Desktop clips. The newer run used NumPy 2.2.6
and ezc3d 1.6.3, while the selected run used NumPy 2.4.4 and ezc3d 1.7.2;
both used MATLAB R2025b Update 5. Model and runtime changed together, so
these observations alone cannot identify the cause of the difference.
The owner export exited naturally with code zero at
\texttt{2026-10-02T08:39:18Z}, improving mean/max position RMS to
16.643/36.865 mm. All current-model clips fully decoded. Owner left foot
mean/max angular errors are 4.517/19.556 degrees and right 9.421/20.299
degrees, with 45 measured samples per foot out of 46. Orientation changes
are mixed: left maximum and right mean worsen slightly. The selected Desktop
package uses the earlier model for the tour and the current model for the
owner, recording each separately. It selects exploratory position tracking,
not anatomical or physical acceptance. Canonical SI XYZ, validity masks,
ordered-label digests and capture rates match exactly under the two recorded
NumPy/ezc3d runtimes. This narrows importer concerns but does not establish
solver equivalence or explain the model difference. A controlled native
model/default comparison was then performed using an immutable private copy of
the earlier model with current source and host runtime. Its tour joint poses
are identical to the current-model run, with the same 13.366/41.385 mm
position metrics. The initial wrapper used a nonexistent field on the returned
report after the solve and exited with code one; a corrected wrapper validated
the actual JSON provenance, reused the verified pose cache and exited naturally
with code zero at \texttt{08:59:42Z}. This is a cache-validated continuation,
not a second fresh solve. Recorded solve-source hashes match those of the
earlier selected execution, but the earlier pose arrays differ. Recorded
identity does not cover all transitive native/runtime inputs; the cause of
this cross-execution difference remains unresolved. The newer model binary
alone does not explain the observed tour fit change under the tested common
conditions. Software tests, a published draft and natural IK export
receipts do not establish contact, anatomical or forward-dynamics acceptance.

\begin{thebibliography}{9}
\bibitem{ezc3d}
pyomeca. \emph{ezc3d: Python C3D reader documentation}.
\url{https://github.com/pyomeca/ezc3d}.
\bibitem{deLeva1996}
P.~de~Leva. Adjustments to Zatsiorsky--Seluyanov's segment inertia parameters.
\emph{Journal of Biomechanics}, 29(9):1223--1230, 1996.
\href{https://doi.org/10.1016/0021-9290(95)00178-6}{doi:10.1016/0021-9290(95)00178-6}.
\bibitem{MathWorksSpherical}
MathWorks. \emph{Spherical Joint: Simscape Multibody Documentation}.
\url{https://www.mathworks.com/help/sm/ref/sphericaljoint.html}.
Quaternion state representation does not remove the adapter's Euler singularity.
\bibitem{Implementation}
UpstreamDrift. \emph{GS3DX Implementation and October 1 Turnover Record}.
Reviewed source commit \texttt{144e81188dd7bb106f81d89b7a7330dd20cce511}.
Nominal ROM source labels are inherited from \texttt{gs3dx\_joint\_rom.m};
they are implementation priors, not newly established clinical limits.
Historical results are attributed to that turnover unless a fresh receipt
identifies a new execution.
\end{thebibliography}
\end{document}
